\begin{mistake}{2026-05-13}{03}

\begin{problemcontent}
当\(x \to 0\)时，\(x - \sin x\cos x\)与\(ax^b\)为等价无穷小，则\(a =\)\_\_\_\_\_\_，\(b =\)\_\_\_\_\_\_。
\end{problemcontent}

\begin{solutioncontent}
利用二倍角公式化简后做泰勒展开。

\(\sin x\cos x = \dfrac{1}{2}\sin 2x\)，故
\[
x - \sin x\cos x = x - \frac{1}{2}\sin 2x
\]

对\(\sin 2x\)做泰勒展开：
\[
\sin 2x = 2x - \frac{(2x)^3}{3!} + O(x^5) = 2x - \frac{4x^3}{3} + O(x^5)
\]

代入：
\[
x - \frac{1}{2}\left(2x - \frac{4x^3}{3}\right) = x - x + \frac{2x^3}{3} = \frac{2}{3}x^3
\]

因此\(x - \sin x\cos x \sim \dfrac{2}{3}x^3\)，即\(a = \dfrac{2}{3}\)，\(b = 3\)。
\end{solutioncontent}

\mistakeknowledge{
\begin{itemize}
  \item \textbf{核心公式/定理}：\(\sin x\cos x = \frac{1}{2}\sin 2x\)；\(\sin u = u - \frac{u^3}{6} + O(u^5)\)
  \item \textbf{易错点}：展开\((2x)^3\)时漏算系数8；忘记外层\(\frac{1}{2}\)的系数
  \item \textbf{关联知识}：等价无穷小的定义——主项系数即为\(a\)，主项幂次即为\(b\)
\end{itemize}}

\end{mistake}
